\section{ТЕХНИЧЕСКОЕ ЗАДАНИЕ}

\subsection{Наименование}

Полное название: Разработка гибридной системы интеллектуального нормоконтроля с интеграцией LLM для семантического анализа технической документации.

Краткое название: Автоматизация процессов контроля технической документации.

\subsection{Назначение}

В процессе ручного контроля проектной документации, 
проверяющий выполняет достаточно рутинную работу, 
для которой характерно возникновениe различных типовых ошибок. 
Поддержание высокого качества результатов, 
на протяжении длительного цикла проверки документа, является трудоемкой задачей. 
Требования для сверок и документы, подлежащие обработке, хранятся децентрализованно, в различных форматах. 
Сам способ хранения и формат может варьироваться для конкретной итерации. 
Понятие централизации экспертных заключений над ошибками размыто. 
Делегирование полномочий между участниками препятствует непрерывному циклу согласований, 
а распределение нагрузки между сотрудниками не всегда прозрачно и эффективно. 
Количество ошибок резко растет при одновременном ведении нескольких проектов. 
Все перечисленные факторы негативно влияют на скорость и стабильность процесса технического контроля. 


Предлагаемая система автоматизированного нормоконтроля 
и лингвистического анализа проектной документации направлена 
на улучшение ключевых метрик за один цикл проверки, 
снижение влияния человеческого фактора, 
а также добавление в каждую итерацию механизма двойной проверки, 
сочетающего автоматизированный и экспертный контроль. 
В рамках системы предусматривается возможность импорта требований, 
что способствует созданию единого источника данных. 
Механизм разграничения по ролям устранит скрытую передачу 
задач и перегрузку команды из-за недобросовестного делегирования. 
В системе будет предусмотрен выбор требований для конкретного проектного 
документа и цикла проверки. Обнаружение лингвистических ошибок позволит 
усилить контроль официально-деловой стилистики. Предусмотренный 
механизм версионирования рабочей документации и хранения файлов в 
MinIO позволит формировать отчеты на основе сравнения версий. 
Вывод предупреждений и пояснений о расхождении с предъявленными 
требованиями добавит интерактивность редактору за счёт использования 
контекстных подсказок. Система строится на базе микросервисной архитектуры 
и реализует гибкий конвейер обработки, где проектный документ будет проходить 
две стадии контроля: синтаксический и семантический. Гибкость конвейера 
будет осуществляться посредством оркестратора Camunda 7. 
Использование стандарта BPMN позволяет просматривать, 
строить и изменять рабочие процессы верификации без необходимости 
пересборки программного кода. BPMS будет доработана специализированными 
блоками действий (actions) для легкой интеграции агентов в рабочий процесс и 
создания Human-in-the-Loop AI подхода. Анализ документа, сравнение версий и 
формирование подсказок предлагается осуществлять посредством self-hosted 
модели в сочетании с RAG архитектурой, которая позволит работать в 
рамках заданных требований. Требования, поступающие от заказчика, 
для создания RAG контекста, будут сегментироваться методом Recursive Chunking 
и индексироваться в векторной базе данных ChromaDB.

\subsection{Основание для разработки}
При разработке используются следующие источники:
\begin{enumerate}[label=\arabic*., nosep, leftmargin=*, labelindent=\parindent]
  \item Кузнецова Е. С., Федюнина Е. А., Орлов И. А., Чукин П. Е. Разработка алгоритма процесса согласования документов в 1С: PLM // Электронные информационные системы. 2022. № 1 (32). С. 48–58. EDN: QWTYLC.
  \item Баданина Н. Д., Судаков В. А., Яшин Н. А. Программный комплекс визуализации и моделирования на основе BPMN нотации // Научная визуализация. 2022. Т. 14. № 3. С. 13–28. DOI: 10.26583/sv.14.3.02. EDN: EQVPDK.
  \item Шкуропадский И. В., Овчинников П. В., Ткачев А. Н. Разработка систем интеллектуальной автоматизации бизнес-процессов на основе технологий машинного обучения // Друкеровский вестник. 2021. № 6 (44). С. 45–56. DOI: 10.17213/2312-6469-2021-6-45-56. EDN: YRJPBQ.
\end{enumerate}

\subsection{Функции}
Предлагаемая система автоматизации должна:
\begin{itemize}[nosep, leftmargin=*, labelindent=\parindent]
  \item Система должна обеспечивать просмотр списка стандартов документов;
  \item Система должна позволять создавать новые стандарты документов;
  \item Система должна обеспечивать редактирование существующих стандартов;
  \item Система должна позволять удалять стандарты документов;
  \item Система должна обеспечивать просмотр списка документов;
  \item Система должна позволять просматривать детальную информацию по выбранному документу;
  \item Система должна обеспечивать сохранение документа в системе;
  \item Система должна позволять сохранять версии документа с возможностью их последующего сравнения;
  \item Система должна обеспечивать выбор конкретной версии документа для просмотра или восстановления;
  \item Система должна позволять загружать документы в систему;
  \item Система должна обеспечивать скачивание документов на локальное устройство пользователя;
  \item Система должна позволять удалять документы из системы;
  \item Система должна обеспечивать запуск автоматической проверки документа на соответствие установленным стандартам;
  \item Система должна выводить список предупреждений и ошибок по результатам проверки;
  \item Система должна позволять выбирать требования для проверки документа из списка доступных проверок;
  \item Система должна обеспечивать сравнение двух версий документа с визуализацией различий;
  \item Система должна позволять выбирать версии для сравнения из списка доступных версий документа;
  \item Система должна отображать изменения между версиями в удобном для восприятия формате;
  \item Система должна обеспечивать просмотр списка всех workflows;
  \item Система должна позволять просматривать детальную информацию по выбранному workflow;
  \item Система должна обеспечивать создание новых workflows;
  \item Система должна позволять редактировать существующие workflows;
  \item Система должна обеспечивать удаление workflows из системы;
\end{itemize}

\subsection{Структура}
Рассматриваемая система автоматизации проверок состоит из следующих компонентов:
\begin{itemize}[nosep, leftmargin=*, labelindent=\parindent]
  \item RAG-service — основной сервис для работы с документами и моделями;
  \item lang-service — автономный сервис лингвистических проверок;
  \item models store — хранилище языковых моделей на базе Ollama сервера;
  \item Minio — объектное хранилище для файлов проектов и правил;
  \item ChromaDB — векторная база данных для хранения эмбеддингов документов;
\end{itemize}

\subsubsection{RAG-service}
Центральный микросервис системы, для оркестрации всех операций с документами. 
Включает document controller для обработки HTTP-запросов, 
document service для бизнес-логики работы с документами, 
model controller и model adapter для взаимодействия с языковыми моделями. 
Содержит model service для управления моделями, 
parser service для разбора документов, 
chunk service для рекурсивного разбиения текста на фрагменты и 
embedding model service для получения векторных представлений текста через nomic-модель. 
MinioStore обеспечивает интеграцию с объектным хранилищем, а компонент version отвечает за версионирование документов.

\subsubsection{lang-service}
Сервис для языковой обработки документов, 
функционирующий независимо от основного RAG-service.

\subsubsection{models store}
Хранилище для языковых моделей на базе Ollama сервера, 
содержит набор обученных моделей.

\subsubsection{Minio}
Объектное хранилище, организованное по стандарту AWS,
для управления процессными файлами. 
Содержит: файлы с правилами проверки документов, проектные файлы пользователей и файлы с результатами сравнения версий документов.

\subsubsection{ChromaDB}
Векторная база данных для хранения эмбеддингов разделённых по смыслу документов, полученных через embedding model service. 
Обеспечивает плотный поиск по документам и извлечение релевантных векторов для формирования контекста.

\subsection{Пользовательский интерфейс}
В ходе работы будет создан интерфейс со следующими элементами:
\begin{itemize}[nosep, leftmargin=*, labelindent=\parindent]
  \item На главной странице системы будет добавлен раздел для управления документами с требованиями. В данном разделе пользователь сможет загружать, просматривать, редактировать и удалять документы, содержащие централизованные требования к проверке проектных документов;
  \item На главной странице системы будет добавлен раздел для управления проектными документами. В данном разделе будут отображаться текущие документы, находящиеся на стадии верификации, а также архив проверенных документов со всеми их версиями, сгруппированными по проектам. При нажатии на документ будет открываться страница с детальной информацией и историей версий;
  \item На странице документа будет добавлен раздел сравнения версий. Пользователь сможет выбрать две версии одного документа из выпадающего списка и запустить сравнение. После обработки запроса система отобразит модальное окно с резюме сравнения, включающим список отличий, добавленных и удаленных элементов между выбранными версиями;
  \item На странице документа будет добавлен раздел автоматической проверки. Пользователь сможет выбрать версию документа для проверки, затем указать требования из списка доступных правил проверки и запустить верификацию. После завершения проверки система отобразит окно с результатами, содержащее список ошибок и предупреждений с указанием проблемных элементов документа. На данном разделе будет доступна кнопка скачивания документа для локального редактирования и кнопка повторной загрузки исправленной версии для следующей итерации проверки;
\end{itemize}

\subsection{Нефункциональные требования}

\subsubsection{Надежность}
В сервисе автоматизации для отказоустойчивой работы нужно предусмотреть следующие аспекты:
\begin{itemize}[nosep, leftmargin=*, labelindent=\parindent]
  \item Должно быть обеспечено еженедельное резервное копирование документов требований на корпоративные хранилища;
  \item Должно быть обеспечено ежедневное резервное копирование проектных документов и всех их версий;
  \item В случае падения модели из сервиса моделей должна оперативно подниматься копия обученной модели данного типа;
  \item В случае падения сервиса для поиска лингвистических ошибок, должен подниматься второй экземпляр сервиса;
\end{itemize}

\subsubsection{Безопасность}
В сервисе автоматизации должна быть предусмотрена ролевая модель, которая поможет разграничить доступ к функционалу.

Контролёр имеет доступ только к документным операциям:
\begin{itemize}[nosep, leftmargin=*, labelindent=\parindent]
  \item Манипуляции с документами требований;
  \item Манипуляции с проектными документами;
  \item Запросы на сравнение версий;
  \item Запросы на проверку проектных документов;
\end{itemize}

Администратор имеет к документным операциям и к операциям цикла:
\begin{itemize}[nosep, leftmargin=*, labelindent=\parindent]
  \item Манипуляции с документами требований;
  \item Манипуляции с проектными документами;
  \item Запросы на сравнение версий;
  \item Запросы на проверку проектных документов;
  \item Манипуляции с процессными схемами;
\end{itemize}

\subsubsection{Условия эксплуатации}
Платформа по верификации документов используется в компаниях с наличием отдела нормоконтроля. Приблизительное количество проверок документов в день 30. 
Система может обеспечивать хранение проектных документов в количестве 1000 штук по 40 страниц. Архив с требованиями к проверке может содержать до 20000 документов.
Приблизительное количество пользователей может варьироватся от 10 до 20 человек.

\subsubsection{Прочие требования}
Все хранилища данной системы должны быть приспособлены к георепликации данных, резервному копированию, архивированию неактуальных документов.

Отклик системы должен позволять делать до 30 проверок в день без потери производительность. 
Система должна позволять делать 10 проверок одновременно без увеличения времени отклика.

Сравнение версий документов не должно длиться более 1 минуты с учётом максимального количества параллельных сравнений, которое достигает 20 штук.

\subsection{Документация}
По окончании работы, будут сформированы следующие документы:
\begin{enumerate}[label=\arabic*., nosep, leftmargin=*, labelindent=\parindent]
  \item Руководство пользователя с описанием интерфейса и основных сценариев использования системы
  \item Руководство администратора с инструкциями по развертыванию и настройке системы
  \item API-документация для всех REST-endpoints сервисов RAG-service и lang-service
  \item Архитектурная документация системы с описанием всех сервисов и их взаимодействия
\end{enumerate}

\subsection{Стадии и этапы разработки}
Система будет разрабатываться по плану:
\begin{enumerate}[label=\arabic*., nosep, leftmargin=*, labelindent=\parindent]
  \item Проектирование архитектуры системы и разработка диаграмм компонентов
  \item Проектирование базы данных и структуры и разметка Minio и ChromaDB
  \item Разработка интерфейсов для RAG-service и lang-service
  \item Развертывание Ollama server и загрузка языковых моделей
  \item Разработка базовых компонентов RAG-service: document controller, document service, MinioStore
  \item Реализация модуля версионирования документов
  \item Разработка parser service для обработки различных форматов документов
  \item Реализация chunk service для разбиения текста
  \item Интеграция embedding model service с nomic-моделью
  \item Настройка ChromaDB и реализация модуля сохранения эмбеддингов
  \item Разработка lang-service для языковой обработки и генерации ответов
  \item Реализация модуля сравнения версий документов
  \item Разработка модуля автоматической проверки документов на соответствие требованиям
  \item Реализация пользовательского интерфейса для работы с документами
  \item Реализация интерфейса управления правилами проверки и стандартами
  \item Функциональное тестирование системы
  \item Нагрузочное тестирование и оптимизация производительности
  \item Проверка в системе антиплагиат, внесение изменений по заключению
  \item Подготовка документации и руководств пользователя
  \item Подготовка презентации, пояснительной записки, текста ВКР и защитной речи
  \item Предзащита и получение допуска к защите
  \item Защита ВКР перед комиссией
  \item Ответы на вопросы комиссии
\end{enumerate}

\subsection{Порядок контроля и приемка}
Этапы:
\begin{enumerate}[label=\arabic*., nosep, leftmargin=*, labelindent=\parindent]
  \item Проведение демонстрации системы научному руководителю от компании
  \item Проведение демонстрации системы научному руководителю от Университета ИТМО
  \item Развёртывание системы локально, предоставление возможности воссоздания работы с документами представителям отдела нормоконтроля
  \item Сбор обратной связи по тестированию системы сотрудниками нормоконтроля
  \item Обработка обратной связи на итоговые метрики улучшения процесса, при переходе от ручных проверок к автоматизированным
  \item Проверка оформления текста ВКР на соответствие стандратам Университета ИТМО и ГОСТ 7.32–2017
  \item Загрузка текста ВКР в систему «Антиплагиат»
  \item Защита ВКР перед комиссией
\end{enumerate}